\documentclass[twoside]{article}
\usepackage{package}
\usepackage[T1]{fontenc}
\usepackage{subfiles}
% Silenciar la advertencia de versalita + negrita inexistente
\WarningFilter{latex}{Some font shapes were not available}
\WarningFilter{latex}{Font shape `T1/lmr/bx/sc'}

% Silenciar la advertencia de caracteres perdidos en el nullfont (para 01_introduccion.tex)
\WarningFilter{latex}{Missing character: There is no}
%Nuevos comandos
\renewcommand{\thefootnote}{\arabic{footnote}}
\raggedbottom



\begin{document}
\begin{center}
    \thispagestyle{empty}
    \large{\textsc{Matrices en MATLAB y Respuesta Transitoria}}\\[0.5cm] 
    \large{Fabrizio Benavidez\footnote[1]{\texttt{fabrizio.benavidez@mi.unc.edu.ar}}}\\[0.5cm]

    \large{\footnotemark[1]\textit{Facultad de Ciencias Exactas, Físicas y Naturales, Ingeniería Aeronáutica}}
    \\[0.5cm]

    \begin{minipage}{0.7\textwidth}
    \hrule
    \vspace{0.2cm}
        \normalsize{\textsc{Resumen}. El presente trabajo aborda la implementación computacional de algoritmos numéricos y el análisis de sistemas de control dinámicos mediante MATLAB y Simulink. En una primera etapa, se desarrollan funciones iterativas utilizando bucles \texttt{for} y \texttt{while} para el cálculo de la sucesión de Fibonacci y la evaluación de límites numéricos en números primos. Posteriormente, se implementa el criterio de estabilidad de Routh-Hurwitz para la evaluación algebraica de polinomios característicos de cuarto orden. Finalmente, se estudia la respuesta en estado estacionario de un sistema de control retroalimentado sujeto a diferentes señales de entrada (escalón, rampa y parábola). Mediante la aplicación del Teorema del Valor Final y la variación de la función de transferencia de retroalimentación $H(s)$, se analiza el impacto del tipo de compensador sobre el error estacionario $e_{ss}$ y la ubicación de los polos en lazo cerrado para determinar la estabilidad del sistema.}\\[0.2cm]
        \normalsize{\textit{Palabras claves}: MATLAB, Simulink, ciclos, estabilidad.\ }
    \vspace{0.2cm}
    \hrule
    \end{minipage}
    
    \vfil

    \includegraphics[width = 0.4\textwidth]{imagenes/Matlab_Logo.png}\\[0.5cm]

    \vfil

    \today\\[0.2cm]
    \large{Semana 5}
\end{center}





\newpage

\section*{Guía 09} 
\subfile{secciones/guia09}

\newpage

\section*{Guía 10} 
\subfile{secciones/guia10}














\newpage
% \thispagestyle{empty}
% \bibliography{biblioteca.bib}
% \bibliographystyle{plain}


\end{document}